\documentclass[11pt]{article}
\usepackage[T1]{fontenc}
\usepackage{lmodern,amsmath,amssymb,amsthm,mathtools,booktabs,microtype}
\usepackage[margin=1.08in]{geometry}
\usepackage[colorlinks=true,linkcolor=blue!60!black,urlcolor=blue!60!black,citecolor=blue!60!black]{hyperref}
\usepackage{xcolor}
\definecolor{navy}{HTML}{203D5C}
\hypersetup{linkcolor=navy,citecolor=navy,urlcolor=navy}
\newtheorem{theorem}{Theorem}[section]
\newtheorem{lemma}[theorem]{Lemma}
\newtheorem{corollary}[theorem]{Corollary}
\newtheorem{proposition}[theorem]{Proposition}
\theoremstyle{remark}\newtheorem{remark}[theorem]{Remark}
\newcommand{\C}{\mathbb C}
\newcommand{\N}{\mathbb N}
\newcommand{\D}{\mathcal D}
\newcommand{\supp}{\operatorname{supp}}
\newcommand{\diag}{\operatorname{diag}}
\newcommand{\norm}[1]{\left\lVert#1\right\rVert}
\newcommand{\abs}[1]{\left|#1\right|}
\newcommand{\eps}{\varepsilon}
\newcommand{\one}{\mathbf 1}
\title{\textbf{Error-correcting prime samples for sparse\\Dirichlet character mixtures}\\[4pt]\large Near-optimal adversarial recovery from polylogarithmically small primes}
\author{Research note for the ProveIt GRH project}
\date{27 November 2026}
\begin{document}
\maketitle
\begin{abstract}
Assuming Dirichlet GRH, we construct a universal set of distinct primes on which every sparse complex mixture of primitive Dirichlet characters of conductor at most $Q$ can be decoded by a single convex optimization problem, even with arbitrary errors and known erasures. For fixed $0<\eps<1$ and sparsity $k\le\abs{\D_Q}\le Q^2$, the number of samples is $O_\eps(k^2\log(2Q))$ and the primes are $O_\eps(k^2\log^2(2Q))$. Recovery is exact whenever the fraction erased plus twice the fraction corrupted is less than $(1-\eps)/(2k)$; an explicit $\ell^1$ stability bound also holds for compressible coefficients and noisy data. Conversely, as soon as the dictionary contains a suitable cyclic character family, no sampling set of primes or decoder can guarantee recovery once this combined budget reaches $1/(2k)$ (up to integer rounding). The obstruction uses only primitive characters of one prime conductor. We derive the analytic sampling estimate from the Dirichlet GRH theorem exported by the vendored Lean development, and extend the merged article's universal prime sampling and single-character recovery results. The finite-dimensional decoding and impossibility arguments are unconditional.
\end{abstract}

\section{Statement and source boundary}
Let $Q\ge3$ be an integer. Let $\D_Q$ consist of all primitive Dirichlet characters of conductors $1\le q\le Q$, with exactly one representative at its primitive conductor and with the character of conductor $1$ included. Extend each character by zero at integers not coprime to its conductor. The number $D=\abs{\D_Q}$ is at most $\sum_{q\le Q}\varphi(q)\le Q^2$. For a coefficient vector $c\in\C^{\D_Q}$ define
\[
 f_c(n)=\sum_{\chi\in\D_Q} c_\chi\chi(n),\qquad
 \sigma_k(c)_1=\min_{\abs{S}\le k}\sum_{\chi\notin S}\abs{c_\chi}.
\]
Thus $\sigma_k(c)_1=0$ precisely when $c$ is $k$-sparse, i.e. it has at most $k$ nonzero entries. For a finite set of distinct primes $P=\{p_1,\dots,p_m\}$ let $A_{j,\chi}=\chi(p_j)$. All vector $\ell^1$ norms below are unnormalized. The measured data are
\[
 y_j=(Ac)_j+b_j+w_j\quad(j\notin E),
\]
where $E$ is a known set of erased coordinates, $b$ has at most $t$ nonzero coordinates outside $E$, and $w$ is arbitrary ordinary noise. Neither the support nor the values of $b$ are known to the decoder.

\begin{theorem}[Universal prime decoder; uses Dirichlet GRH]\label{thm:main}
Fix $0<\eps<1$ and $1\le k\le D$. There is a set $P$ of distinct primes, depending only on $Q,k,\eps$, with
\begin{align*}
 m&\asymp_\eps k^2\log(2Q),&
 \max P&\ll_\eps k^2\log^2(2Q).
\end{align*}
Let $e=\abs E$ and put $\Delta=1-\eps-2k(e+2t)/m$. If $\Delta>0$, every minimizer
\begin{equation}\label{eq:decoder}
 \widehat c\in\arg\min_{z\in\C^{\D_Q}}
 \left(\norm{z}_1+\frac{2k}{m}\sum_{j\notin E}\abs{y_j-(Az)_j}\right)
\end{equation}
satisfies
\begin{equation}\label{eq:stability}
 \norm{\widehat c-c}_1\le
 \frac{2\sigma_k(c)_1+(4k/m)\sum_{j\notin E}\abs{w_j}}{\Delta}.
\end{equation}
In particular, exact $k$-sparse recovery holds with no ordinary noise for every error pattern satisfying $e+2t<(1-\eps)m/(2k)$. The decoder does not need to know $t$, the corrupted coordinates, or the nonzero coefficients.
\end{theorem}

\begin{theorem}[Universal obstruction; unconditional given the modulus]\label{thm:lower}
Suppose there exists a prime $q\le Q$ such that $q\equiv1\pmod{4k}$. For every list of $m\ge1$ primes, even with repetitions and no size restriction, there are two different exactly $k$-sparse vectors in $\C^{\D_Q}$ whose measurements differ in at most $\lfloor m/(2k)\rfloor$ positions. Thus if integer budgets satisfy
\begin{equation}\label{eq:lower}
 e+2t\ge\left\lfloor\frac{m}{2k}\right\rfloor,
\end{equation}
no decoder can guarantee exact recovery of every $k$-sparse vector in the presence of up to $e$ known erasures and up to $t$ unknown corruptions. Under Dirichlet GRH the modulus condition holds for all $Q\ge C k^2\log^2(2k)$, with an absolute constant $C$.
\end{theorem}

For example, with $k=2$ one can recover after nearly $25\%$ erasures, or nearly $12.5\%$ arbitrary errors, or any combination on the indicated trade-off line. The claim is about the worst case over unknown mixtures, not a statistical model for errors. Theorems \ref{thm:main} and \ref{thm:lower} give a sharp limiting fraction, not a claim of optimal sample count or sharp finite-$m$ decoding.

\paragraph{Input from the repository.}
The two supplied paths were accessed through GitHub. The vendored readable edition of \texttt{grh-lean} is identified by its README as upstream commit
\begin{center}\small\texttt{72fead102f663d5c7a407964a13dc0cf9d977207}.\end{center}
In \texttt{grh/src/Statements.lean}, the propositions \texttt{DirichletLNonvanishingRightOfHalf} and \texttt{DirichletLZerosOnLine} quantify over all nonzero natural moduli and all complex Dirichlet characters. In \texttt{grh/src/Main.lean}, the corresponding theorems are \path{LFunction_ne_zero_of_half_lt_re} and \path{LFunction_eq_zero_re}; the theorem \texttt{riemannHypothesis} has Mathlib's standard type. The file \texttt{New/Main.lean} uses the functional equation to pass from the right-half-plane result to the critical-line statement. The preceding analytic mechanism, described in \path{New/Reduction/Zeros/Criterion.lean}, uses a smoothed prime-power remainder and a one-sided Landau argument. Here we reverse the direction: the resulting zero theorem supplies a uniform square-root estimate for a different, compactly supported prime statistic. No simplicity, linear independence, or nonabelian GRH is used.

The upstream README reports only \texttt{propext}, \texttt{Classical.choice}, and \texttt{Quot.sound} as axioms of the five theorems; Sutherland's independent verification page reports checks of the shorter editions \cite{lean,sutherland}. The vendor is a separate Lake project (Lean 4.35.0, pinned Mathlib); this note did not rebuild or audit its full proof. The merged article's README explicitly distinguishes results using Dirichlet GRH from unconditional results and reductions. We retain that distinction.

\paragraph{Relation to earlier reports.}
The merged article \cite{article}, Part III (files \texttt{dir.tex} and \texttt{rec.tex}), already contains a universal small-prime list for bounded character bias (label \texttt{dir:lem:oneset}), squarefree product and partition theorems, and recovery of a single primitive character via a short dictionary of primes (label \texttt{rec:cor:dict}). The latter includes corruption tolerance for individual characters; the article also treats corrupted zero-residue data (label \texttt{rec:thm:stable}). Our contribution is the quantitative sparsity-dependent Gram estimate at bias $O(1/k)$, its combination with a convex decoder for arbitrary mixed-conductor complex coefficients, and the sharp matching error/erasure obstruction. The small-prime estimates themselves and uncertainty ideas are classical \cite{mv,tao,ds,jmv}; no novelty is claimed for them individually. Inspection was targeted to the pertinent sections of the 723-page merged article, not a full novelty audit of its 914 listed results or all mathematical literature. The specific combined theorem is proposed as a new result, not certified as historically unprecedented.

\section{Analytic estimate: small primes almost orthogonal}\label{sec:analytic}
Fix once and for all a nonnegative $C^\infty$ function $W$ supported in $[1,2]$, positive in $(1,2)$, with $\int W(u)\,du=1$. Constants may depend on this fixed choice. For example, normalize $\exp(-1/[(u-1)(2-u)])$ on $(1,2)$ and extend by zero. Put
\[
 w_p(x)=(\log p)W(p/x),\qquad H(x)=\sum_p w_p(x).
\]
\begin{lemma}[Smoothed prime estimate; uses Dirichlet GRH]\label{lem:explicit}
For $x\ge2$, every modulus $d\ge1$ and every character $\psi$ modulo $d$,
\begin{equation}\label{eq:prime}
 \sum_p (\log p)W(p/x)\psi(p)
 =\one_{\psi\text{ principal}}x+O_W(\sqrt x\log(2d)).
\end{equation}
The implied constant is independent of $d,\psi,x$. In particular $H(x)=x+O_W(\sqrt x)$.
\end{lemma}
\begin{proof}
Here are the details of the standard smoothed explicit formula, so the uniformity needed later is visible. First take a primitive character $\psi$ of conductor $f$ (including the zeta case $f=1$). The Mellin transform $\widetilde W(s)=\int W(u)u^{s-1}\,du$ is entire and, by integration by parts, decays faster than any power of $\abs{\Im s}$ on fixed vertical strips. Mellin inversion on $\Re s=2$ gives
\[
 \sum_n\Lambda(n)\psi(n)W(n/x)=\frac1{2\pi i}\int_{(2)}
 -\frac{L'}{L}(s,\psi)x^s\widetilde W(s)\,ds.
\]
Shift to $\Re s=1/4$. The only enclosed poles are the pole at $1$ in the zeta case and the nontrivial zeros $\rho$ of $L(s,\psi)$, counted with multiplicity. The residue theorem yields
\begin{equation}\label{eq:explicit}
 \one_{f=1}x-\sum_\rho x^\rho\widetilde W(\rho)+O_W(x^{1/4}\log(2f)).
\end{equation}
For completeness, the standard local zero count is $O(\log(2f(2+\abs T)))$ for ordinates in $[T,T+1]$. The local partial-fraction formula for $L'/L$ gives a sum of reciprocals of distances to zeros within distance $1$ in ordinate, plus $O(\log(2f(2+\abs t)))$ and a pole term in the zeta case. These facts follow from the functional equation and the usual zero-count/logarithmic-derivative estimates (see \cite[Chapter 13]{mv} and \cite[Exercise 45]{tao}). Under GRH the distance in real part from $1/4$ to all nontrivial zeros is $1/4$, so $L'/L$ is $O(\log(2f(2+\abs t)))$ on that line, including bounded $t$. Horizontal contours can be chosen in successive unit height intervals at distance $\gg 1/\log(2f(2+\abs T))$ from the ordinates of nearby zeros. On these contours $L'/L$ has at most logarithmic-squared growth, and rapid Mellin decay makes their integrals vanish as $\abs T\to\infty$. This justifies \eqref{eq:explicit} and the uniform error there.

GRH and the local zero count imply
\[
 \sum_\rho\abs{\widetilde W(\rho)}
 \ll_W\sum_{j\in\mathbb Z}\frac{\log(2f(2+\abs j))}{(1+\abs j)^3}
 \ll_W\log(2f).
\]
Thus the zero term is $O_W(\sqrt x\log(2f))$. To remove prime powers, use the elementary Chebyshev bound $\vartheta(y)=\sum_{p\le y}\log p\ll y$:
\[
 \sum_{a\ge2}\vartheta((2x)^{1/a})
 \ll \sqrt x+x^{1/3}\log(2x)\ll\sqrt x,
\]
where only $a\le\log(2x)/\log2$ occur. Multiplying by $\norm W_\infty$ is harmless. Finally for a character modulo $d$ induced from $\psi^*$ modulo $f\mid d$, the prime sum changes only at primes dividing $d$, by at most $\norm W_\infty\log d$. The inducing character has $f=1$ precisely when the character modulo $d$ is principal. This proves the assertion.
\end{proof}

\begin{lemma}[Uniform weighted Gram matrix]\label{lem:weighted}
Write
\[
 G^w_{\chi,\psi}(x)=H(x)^{-1}\sum_p w_p(x)\overline{\chi(p)}\psi(p).
\]
For all $x$ larger than an absolute threshold (depending on the fixed $W$),
\begin{equation}\label{eq:weighted}
 \max_{\chi,\psi\in\D_Q}\abs{G^w_{\chi,\psi}(x)-\one_{\chi=\psi}}
 \ll_W\frac{\log(2Q)}{\sqrt x}.
\end{equation}
\end{lemma}
\begin{proof}
For conductors $q,r$, inflate both characters to $d=\operatorname{lcm}(q,r)\le Q^2$ and take their quotient on units, extended by zero. On all primes this product agrees with $\overline{\chi(p)}\psi(p)$, including primes dividing $d$. The quotient is principal if and only if the inflated characters agree, which by uniqueness of the inducing primitive character is equivalent to $\chi=\psi$ in our dictionary. Apply Lemma \ref{lem:explicit}, divide by $H(x)\asymp x$, and use $\log(2d)\le2\log(2Q)$. In particular, we did not presume all the sampled primes to be coprime to all possible conductors: the diagonal estimate includes the missing values automatically.
\end{proof}

\section{A universal unweighted sample of distinct primes}\label{sec:sampling}
\begin{theorem}[Quantitative prime coherence; uses Dirichlet GRH]\label{thm:gram}
For $Q\ge3$ and $0<\mu\le1/2$, there exists a set of distinct primes $P$ such that
\begin{align}
 m=\abs P&\asymp \mu^{-2}\log(2Q),\label{eq:m}\\
 \max P&\ll \mu^{-2}\log^2(2Q/\mu),\label{eq:x}\\
 \max_{\chi,\psi\in\D_Q}\left|\frac{1}{m}\sum_{p\in P}\overline{\chi(p)}\psi(p)-\one_{\chi=\psi}\right|&\le\mu.\label{eq:gram}
\end{align}
The constants are absolute once $W$ is fixed, including as $\mu\downarrow0$.
\end{theorem}
\begin{proof}
Let $L=\log(2Q/\mu)$, $x=K\mu^{-2}L^2$ and $M=B\mu^{-2}\log(2Q)$. Take $B$ a sufficiently large absolute constant and then take $K$ sufficiently large in terms of $B,W$. Lemma \ref{lem:weighted} gives maximum Gram deviation at most $\mu/8$ and $H(x)\ge x/2$. Independently select each prime of $[x,2x]$ with probability
\[
 \pi_p=Mw_p(x)/H(x).
\]
These probabilities are valid: $\max_p\pi_p\ll_W (B/K)\log(2Q)\log(2x)/L^2$. Since $\log(2x)\le \log(2K)+2\log(1/\mu)+2\log L\ll \log(2K)+L$, and $L\ge\log 6$, increasing $K$ makes this at most $1$. There are no repetitions in this Bernoulli selection.

Let $B_p$ be the indicator of inclusion and $m=\sum B_p$. Then $\mathbb E m=M$. For each ordered pair put
\[
 Z_{\chi,\psi}=\sum_p B_p\bigl(\overline{\chi(p)}\psi(p)-\one_{\chi=\psi}\bigr).
\]
Each summand's coefficient has modulus at most $1$: off-diagonal this is the character product; on the diagonal it is $0$ or $-1$. Thus the real and imaginary parts have total variance at most $\sum_p\pi_p=M$, and each centered real summand is bounded in absolute value by $1$. Also $\abs{\mathbb E Z_{\chi,\psi}}\le\mu M/8$. Bernstein's inequality gives
\[
 \Pr\bigl(\abs{Z_{\chi,\psi}-\mathbb E Z_{\chi,\psi}}>\mu M/4\bigr)
 \le 4\exp(-c\mu^2M)
\]
for an absolute $c>0$ (apply the real estimate at threshold $\mu M/8$ to both parts). Similarly $\Pr(m\notin[M/2,3M/2])\le2e^{-c'M}$. There are at most $Q^4$ pairs. The union of the bad events has probability less than one for sufficiently large $B$, since $\mu^2M=B\log(2Q)$. For a successful outcome, $m\ge M/2$ and
\[
 \frac{\abs{Z_{\chi,\psi}}}{m}\le\frac{3\mu M/8}{M/2}=\frac{3\mu}{4}\le\mu.
\]
The claimed size and prime bound follow from $m\asymp M$ and $p\le2x$.
\end{proof}

The result strengthens the fixed-bias selection used in the merged article in the specific direction needed here: both the bias and the sample size vary uniformly with $k$. A sampled list can be tested by checking the Gram matrix; there is no need to prove GRH separately for each sample. Since character values are roots of unity or zero, a candidate list has an exact algebraic Gram certificate (see Section \ref{sec:implementation}).

\section{The decoder: a finite-dimensional theorem}\label{sec:decoder}
We isolate an unconditional statement, applicable to any bounded complex dictionary. For a matrix $A\in\C^{m\times D}$ let
\[
 \mu(A)=\max_{i,j}\abs{(A^*A/m-I)_{ij}}.
\]
The diagonal is included, which matters for characters vanishing at ramified primes.

\begin{theorem}[Robust convex decoding]\label{thm:convex}
Suppose $\abs{A_{ji}}\le1$ and $\mu(A)\le\mu$. Let $E$ and the data be as in Section 1 and let $e=\abs E$. For $1\le k\le D$ define
\[
 d_0=1-2k\mu-2k\frac{e+2t}{m}.
\]
If $d_0>0$, a minimizer of \eqref{eq:decoder} exists, and every minimizer satisfies \eqref{eq:stability} with $d_0$ in place of $\Delta$. For $\sigma_k(c)_1=0$ and $w=0$ the minimizer is unique and equals $c$.
\end{theorem}
\begin{proof}
The objective is continuous and coercive because of the term $\norm{z}_1$, so at least one minimizer exists. Let $U=E^c$ be the retained row indices, $T\subseteq U$ the support of $b$ with $\abs T\le t$, and let $S$ be a set of $k$ largest coordinates of $c$ in modulus. For any perturbation $h\in\C^D$,
\begin{equation}\label{eq:coord}
 \norm h_\infty
 \le \norm{A^*Ah/m}_\infty+\mu\norm h_1
 \le \frac{\norm{(Ah)_U}_1}{m}+(\mu+e/m)\norm h_1.
\end{equation}
The first inequality is entrywise control of the Gram error. For the second, $\abs{A_{ji}}\le1$ gives $\norm{A^*v}_\infty\le\norm v_1$ and $\norm{Ah}_\infty\le\norm h_1$, so the erased rows cost at most $e\norm h_1$. Multiplying \eqref{eq:coord} by $k$ gives
\begin{equation}\label{eq:support}
 \norm{h_S}_1\le \frac{k}{m}\norm{(Ah)_U}_1+k(\mu+e/m)\norm h_1.
\end{equation}
The coefficient penalty changes by at least
\begin{equation}\label{eq:coeff}
 \norm{c+h}_1-\norm c_1\ge\norm h_1-2\norm{h_S}_1-2\sigma_k(c)_1.
\end{equation}
This follows by the triangle inequality separately on $S$ and $S^c$. For the residual penalty, first put $w=0$. On $U\setminus T$ its change is $\abs{(Ah)_j}$, and on $T$ it is at least $-\abs{(Ah)_j}$. Inserting arbitrary $w$ costs at most $2\norm{w_U}_1$ by two triangle inequalities. Consequently
\begin{equation}\label{eq:residual}
 \norm{b_U+w_U-(Ah)_U}_1-\norm{b_U+w_U}_1
 \ge \norm{(Ah)_U}_1-2t\norm h_1-2\norm{w_U}_1.
\end{equation}
Denote the objective of \eqref{eq:decoder} by $J$. Combining \eqref{eq:support}--\eqref{eq:residual} gives the key cancellation:
\begin{align*}
 J(c+h)-J(c)
 &\ge \norm h_1-2\norm{h_S}_1-2\sigma_k(c)_1\\
 &\quad+\frac{2k}{m}\bigl(\norm{(Ah)_U}_1-2t\norm h_1-2\norm{w_U}_1\bigr)\\
 &\ge d_0\norm h_1-2\sigma_k(c)_1-\frac{4k}{m}\norm{w_U}_1.
\end{align*}
The two terms involving $\norm{(Ah)_U}_1$ cancel exactly at the penalty weight $2k/m$. At a minimizer take $h=\widehat c-c$ and use $J(c+h)-J(c)\le0$. Division by $d_0>0$ proves the bound, and the noiseless sparse case gives $h=0$.
\end{proof}

\begin{proof}[Proof of Theorem \ref{thm:main}]
Use Theorem \ref{thm:gram} with $\mu=\eps/(2k)\le1/2$. Since $k\le D\le Q^2$ and $\eps$ is fixed, $\log(2Q/\mu)\ll_\eps\log(2Q)$. The two size bounds become those of Theorem \ref{thm:main}, while Theorem \ref{thm:convex} has $d_0\ge\Delta$. All choices are independent of $c,E,b,w$. This proves the theorem.
\end{proof}

\begin{corollary}[Minimum distance of sparse prime measurements]\label{cor:dist}
For the sample in Theorem \ref{thm:gram}, if $c\ne c'$ are both $k$-sparse then
\[
 \abs{\supp A(c-c')}\ge m\left(\frac1{2k}-\mu\right).
\]
\end{corollary}
\begin{proof}
For $h=c-c'$ with $s=\abs{\supp h}\le2k$, the Gram estimate gives
\[
 m^{-1}\norm{Ah}_2^2\ge\norm h_2^2-\mu\norm h_1^2\ge(1-s\mu)\norm h_2^2.
\]
Also $\norm{Ah}_\infty^2\le\norm h_1^2\le s\norm h_2^2$. Writing $r=\abs{\supp Ah}$, we obtain $r/m\ge 1/s-\mu\ge1/(2k)-\mu$, as claimed.
\end{proof}
The distance bound alone gives a possibly nonconvex nearest-neighbor decoder. Theorem \ref{thm:convex} shows that a tractable convex objective attains the same guaranteed corruption fraction, even with noisy and compressible inputs. This is the main finite-dimensional improvement over merely separating pairs of characters.

\section{Sharpness by a cyclic packet of primitive characters}
\begin{proof}[Proof of Theorem \ref{thm:lower}]
Take a generator $\theta$ of the cyclic character group modulo the prime $q$, and let $\omega=\theta^{(q-1)/(2k)}$, of order $2k$. The $2k$ characters
\[
 \chi_j=\theta\omega^j,\qquad 0\le j<2k,
\]
are all distinct and nonprincipal: $\theta$ does not lie in the subgroup generated by $\omega$, because $q-1\ge4k$. Every nonprincipal character modulo a prime is primitive. All these characters therefore belong to $\D_Q$ and all vanish at $q$.

For each $2k$th root of unity $v$, form the mixture
\[
 F_v(n)=\sum_{j=0}^{2k-1} v^{-j}\chi_j(n).
\]
On integers coprime to $q$, the geometric sum gives
\begin{equation}\label{eq:comb}
 F_v(n)=\begin{cases}2k\,\theta(n)&\omega(n)=v,\\0&\omega(n)\ne v.\end{cases}
\end{equation}
At multiples of $q$ all summands vanish. Among the $2k$ disjoint classes of sample coordinates defined by $\omega(p)=v$, choose one of minimum size $r\le\lfloor m/(2k)\rfloor$. Split the sum into two $k$-term halves with opposite signs: set $c_{\chi_j}=v^{-j}$ for $j<k$ and $c'_{\chi_j}=-v^{-j}$ for $k\le j<2k$, with other entries zero. Then $A(c-c')=F_v|_P$. The vectors are different, exactly $k$-sparse, and differ in their observed values at precisely $r$ coordinates.

If $e+2t\ge r$, erase up to $e$ of those positions so that at most $2t$ remain, then partition the remaining differing positions into two groups of size at most $t$. On one group use the observations of $c$, on the other those of $c'$, and on all positions where the observations agree use their common value. This single received word is consistent with either coefficient vector, with the same known erasures and at most $t$ arbitrary errors in each case. Hence no rule can guarantee recovery. In particular \eqref{eq:lower} suffices, regardless of the distribution of the sampled primes.

To obtain a suitable modulus under GRH, apply Lemma \ref{lem:explicit} and character orthogonality modulo $4k$:
\[
 \sum_{p\equiv1\ (4k)}w_p(X)=\frac{X}{\varphi(4k)}+O_W(\sqrt X\log(8k)).
\]
Take $X=K_0 k^2\log^2(2k)$ with $K_0$ large enough. The main term dominates since $\varphi(4k)\le4k$ and $\log(8k)\ll\log(2k)$. Positivity supplies a prime $q\le2X$ congruent to $1$ modulo $4k$. This proves the final assertion.
\end{proof}

\paragraph{Quantifying the gap.}
For a sample with Gram bound $\mu$ the minimum Hamming distance between two $k$-sparse measurement vectors is at least $\lceil m(1/(2k)-\mu)\rceil$ by Corollary \ref{cor:dist}. In the range of Theorem \ref{thm:lower} it is at most $\lfloor m/(2k)\rfloor$. Sending $\eps=2k\mu$ to zero in Theorem \ref{thm:main} proves sharpness of the fractional threshold in the sense of arbitrarily small relative loss. The theorem does not assert that the sample-size constant remains bounded as $\eps\to0$, nor that the distance upper bound holds for dictionaries without the stipulated cyclic packet. The impossibility bound concerns arbitrary complex coefficients, not the smaller class of nonnegative integral multiplicities.

\section{Arithmetic consequences and scope}
\begin{corollary}[Recovering sparse packets of $L$-functions; uses Dirichlet GRH]\label{cor:L}
For $\Re s>1$, let
\[
 D_c(s)=-\sum_{\chi\in\D_Q}c_\chi\frac{L'}{L}(s,\chi).
\]
If at most $k$ of the $c_\chi$ are nonzero, its unknown coefficients are recovered from the normalized prime Dirichlet-series coefficients at the sample of Theorem \ref{thm:main}, with precisely the same adversarial guarantees. For $c_\chi$ nonnegative integers, this recovers the primitive factors and their multiplicities in the product $\prod_\chi L(s,\chi)^{c_\chi}$; the product need not be supplied elsewhere.
\end{corollary}
\begin{proof}
The Euler product and absolute convergence yield $-L'/L(s,\chi)=\sum_n\Lambda(n)\chi(n)n^{-s}$. The coefficient at $p$ in $D_c$ is $(\log p)f_c(p)$, so dividing the supplied coefficient by $\log p$ gives exactly the data model. Theorem \ref{thm:main} applies. If the exponents are nonnegative integers, the displayed logarithmic derivative is the negative logarithmic derivative of the product, and recovery of the vector is recovery of the factors with multiplicities.
\end{proof}
If the $c_\chi$ are integers and the right side of \eqref{eq:stability} is less than $1/2$, rounding the real parts of all recovered coefficients to the nearest integer yields exact recovery, even with ordinary noise. For arbitrary complex coefficients the theorem yields an explicit error bound, rather than exact support recovery for coefficients smaller than that error.

The primitive-conductor convention is essential. If two imprimitive characters are induced by the same primitive character, they agree at every prime not dividing their moduli. Such Euler-factor changes cannot in general be recovered from samples avoiding those divisors. We do not claim uniqueness for a dictionary containing these aliases. Nor is the theorem a result about prime evaluations of arbitrary automorphic representations: the mixed-conductor Gram argument depends on products of Dirichlet characters again being Dirichlet characters.

\section{Reproducibility, algorithm and formalization}\label{sec:implementation}
The accompanying \texttt{check.py} enumerates primitive characters by the Chinese remainder theorem and the unit groups of prime powers, verifies their total number independently via $\sum_{d\mid q}\mu_{\rm Mob}(q/d)\varphi(d)$, selects distinct primes with a fixed random seed, and measures the full complex Gram deviation. It then implements the real subdictionary case of the decoder as a linear program, and constructs the explicit sharpness witness. The experiment uses the compactly supported polynomial bump $30(u-1)^2(2-u)^2$ on $[1,2]$; the proof uses a fixed smooth bump, and the experiment tests the resulting finite matrix directly. For complex unknowns, introduce real epigraph variables $a_i,r_j$ with $a_i\ge\abs{z_i}$ and $r_j\ge\abs{y_j-(Az)_j}$; minimizing $\sum_i a_i+(2k/m)\sum_{j\notin E} r_j$ is a second-order cone program in real coordinates.

\begin{center}
\begin{tabular}{@{}lr@{}}
\toprule
Numerical illustration & Value\\\midrule
Maximum conductor $Q$ / sparsity $k$ & $31$ / $2$\\
Prime interval & $[40000,80000]$\\
Distinct samples $m$ / primitive dictionary size & $1616$ / $197$\\
Measured maximum Gram deviation & $0.0607388318$\\
Rigorously certified Gram deviation bound & $0.061$\\
Known erasures / unknown corruptions & $61$ / $91$\\
Noiseless coefficient $\ell^1$ error (real LP) & $3.34\times10^{-13}$\\
Noisy coefficient $\ell^1$ error (real LP) & $0.00348748$\\
Stability upper bound using measured deviation & $0.195109$\\
Sharpness packet support, modulo $17$ & $392$\\
Errors per candidate in constructed collision & $196$\\\bottomrule
\end{tabular}
\end{center}
The LP experiment is numerical and only the real subdictionary decoder is executed. However, the attached \texttt{certify\_gram.py} independently certifies the full complex Gram bound $\mu\le61/1000$ for all 197 primitive characters using rational and integer arithmetic. It constructs exact rational phases by the Chinese remainder theorem, brackets $\pi$ by Machin's identity and alternating arctangent series, and bounds the degree-40 Taylor remainders for sine and cosine by $7^{41}/41!$. Rounding to integers at scale $S=10^6$ makes the error in each complex entry at most $2/S$. Each Gram entry then has error at most $4/S+4/S^2$. Integer matrix products (with a checked no-overflow bound) and exact squared comparisons certify the target. In particular, $1-4(61/1000)-4(61+2\cdot91)/1616=7803/50500>0$ is a rigorous recovery margin for the full dictionary in this finite instance. The LP's stability bound in the table uses the measured deviation; the exact certified bound is slightly weaker. The file \texttt{sample\_primes.txt} records the exact selected list. For arbitrary finite instances one could instead compute Gram entries in the cyclotomic field generated by all character values, and compare their squared absolute values with a rational target $\mu^2$ in the real algebraic number field. No numerical experiment is needed for the mathematical theorem. A deterministic (not necessarily efficient) sample-finding algorithm is to enumerate finite prime sets with size in the proved range and accept the first passing this algebraic test. The existence theorem guarantees termination if the constants are chosen as in its proof; without explicit constants, one may dovetail over increasing prime bounds and sample sizes. The latter does not by itself certify the asymptotic size bounds for the first accepted sample.

\paragraph{Lean integration.}
The proof is organized around two independent interfaces. First, formalize the smoothed estimate of Lemma \ref{lem:explicit} using the exported Dirichlet GRH proposition, Mellin inversion, the functional equation, and a uniform local zero count; prove the Gram and Bernoulli selection lemmas. Second, formalize the finite-dimensional inequalities \eqref{eq:coord}--\eqref{eq:residual} for a bounded matrix on finite index types and derive the decoder theorem for any minimizer; the cyclic packet obstruction uses only finite character groups and a geometric sum. The attached \texttt{formalization.md} gives a dependency plan and exact mathematical interfaces, not a compiled Lean file. The environment used for this note has no Lean executable, so neither this theorem nor its connection to the vendored package has been machine-checked here. No new axioms about L-functions have been introduced beyond Dirichlet GRH as stated.

\paragraph{Research questions.}
The current sample bound comes from a union bound on pairwise Gram entries; a sparsity-sensitive random matrix method might reduce it. The sharp adversarial fraction comes from a subgroup packet of characters and cannot be improved for the full dictionary, even with larger primes. For restricted dictionaries without such packets, the best threshold may instead be controlled by sharper uncertainty inequalities. These are directions, not proved improvements.

\begin{thebibliography}{9}
\bibitem{lean} L. Sernau et al., \emph{GRH for Dirichlet L-functions, in Lean 4}, readable edition, upstream commit \texttt{72fead102f663d5c7a407964a13dc0cf9d977207} (25 Nov. 2026). \url{https://github.com/lsernau/grh-lean}.
\bibitem{sutherland} A. V. Sutherland, \emph{Shorter GRH Lean proof: verification files}, \url{https://math.mit.edu/~drew/shortgrh/} (accessed 27 Nov. 2026).
\bibitem{article} ProveIt Contributors, \emph{Beyond the Riemann hypothesis for Dirichlet L-functions: consequences, transfers and limits}, merged article, 26 Nov. 2026. Repository directory \href{https://github.com/VladimirReshetnikov/ProveIt/tree/main/Analysis/GeneralizedRiemannHypothesis/Article}{Analysis/GeneralizedRiemannHypothesis/Article}, especially sections \texttt{src.tex}, \texttt{dir.tex}, \texttt{rec.tex}. See \texttt{source\_audit.md} for blob identifiers.
\bibitem{mv} H. L. Montgomery and R. C. Vaughan, \emph{Multiplicative Number Theory I: Classical Theory}, Cambridge University Press, 2007, Chapter 13.
\bibitem{tao} T. Tao, \emph{254A, Notes 2: Complex-analytic multiplicative number theory}, 9 Dec. 2014, \href{https://terrytao.wordpress.com/2014/12/09/254a-notes-2-complex-analytic-multiplicative-number-theory/}{online notes}.
\bibitem{ds} D. L. Donoho and P. B. Stark, Uncertainty principles and signal recovery, \emph{SIAM J. Appl. Math.} 49 (1989), 906--931. \href{https://doi.org/10.1137/0149053}{doi:10.1137/0149053}.
\bibitem{jmv} D. Jao, S. D. Miller and R. Venkatesan, Expander graphs based on GRH with an application to elliptic curve cryptography, \emph{J. Number Theory} 129 (2009), 1491--1504.
\end{thebibliography}
\end{document}
